% GENERATED FILE. Edit canonical lessons, then run npm run generate:books.
% canonical-source-hash: fnv1a64:8da3570878f891fa
% canonical-lessons: MR-C253-valne, MR-C253-anand-ghene, MR-C253-kalpana-karne, MR-C253-kalavne, MR-C253-takrar-karne

\chapter{To turn, To enjoy, To imagine, To inform, To complain}
\label{ch:mr-to-turn-to-enjoy-to-imagine-to-inform-to-complain}
\hlchaptermodality{253}

\begin{chapteropening}
\textbf{By the end of this chapter:} \emph{I can say to turn, to enjoy, to imagine, to inform and to complain.}

Complete the last lesson of chapter 253: I can say to turn, to enjoy, to imagine, to inform and to complain.
\end{chapteropening}

\section[\texorpdfstring{\mr{वळणे} (vaḷṇe)}{वळणे (vaḷṇe)}]{\mr{वळणे} (vaḷṇe) — to turn}

\label{lesson:MR-C253-valne}



\begin{quote}
Before the new one: say the Marathi for to smoke, then the Marathi for to go with.
\end{quote}

\subsection*{You'll want to know: \mr{वळणे}}
\textbf{\mr{वळणे}} — \emph{vaḷṇe} — \textquotedblleft{}to turn\textquotedblright{}.

\begin{grammarlens}[title={What you've built}]
One more word for this chapter.
\end{grammarlens}

\subsection*{Your turn}
\begin{itemize}
\raggedright
  \item \emph{Say it:} \emph{vaḷṇe}
  \item \emph{Say it:} \emph{vaḷṇe}, once more
  \item \emph{Say it:} say \emph{vaḷṇe}
  \item \emph{Recall:} say the Marathi for to smoke, then the Marathi for to go with, then say \emph{vaḷṇe} again
\end{itemize}

\begin{tcolorbox}[breakable,colback=teal!4,colframe=teal!35!black,title={Before you move on}]
What does \mr{वळणे} mean? (\textquotedblleft{}To turn\textquotedblright{}.) Say it once more.
\end{tcolorbox}
\section[\texorpdfstring{\mr{आनंद} \mr{घेणे} (ānand gheṇe)}{आनंद घेणे (ānand gheṇe)}]{\mr{आनंद} \mr{घेणे} (ānand gheṇe) — to enjoy}

\label{lesson:MR-C253-anand-ghene}



\begin{quote}
Before the new one: say the Marathi for to go with, then the Marathi for to turn.
\end{quote}

\subsection*{You'll want to know: \mr{आनंद} \mr{घेणे}}
\textbf{\mr{आनंद} \mr{घेणे}} — \emph{ānand gheṇe} — \textquotedblleft{}to enjoy\textquotedblright{}.

\begin{grammarlens}[title={What you've built}]
One more word for this chapter.
\end{grammarlens}

\subsection*{Your turn}
\begin{itemize}
\raggedright
  \item \emph{Say it:} \emph{ānand gheṇe}
  \item \emph{Say it:} \emph{ānand gheṇe}, once more
  \item \emph{Say it:} say \emph{ānand gheṇe}
  \item \emph{Recall:} say the Marathi for to go with, then the Marathi for to turn, then say \emph{ānand gheṇe} again
\end{itemize}

\begin{tcolorbox}[breakable,colback=teal!4,colframe=teal!35!black,title={Before you move on}]
What does \mr{आनंद} \mr{घेणे} mean? (\textquotedblleft{}To enjoy\textquotedblright{}.) Say it once more.
\end{tcolorbox}
\section[\texorpdfstring{\mr{कल्पना} \mr{करणे} (kalpanā karṇe)}{कल्पना करणे (kalpanā karṇe)}]{\mr{कल्पना} \mr{करणे} (kalpanā karṇe) — to imagine}

\label{lesson:MR-C253-kalpana-karne}



\begin{quote}
Before the new one: say the Marathi for to turn, then the Marathi for to enjoy.
\end{quote}

\subsection*{You'll want to know: \mr{कल्पना} \mr{करणे}}
\textbf{\mr{कल्पना} \mr{करणे}} — \emph{kalpanā karṇe} — \textquotedblleft{}to imagine\textquotedblright{}.

\begin{grammarlens}[title={What you've built}]
One more word for this chapter.
\end{grammarlens}

\subsection*{Your turn}
\begin{itemize}
\raggedright
  \item \emph{Say it:} \emph{kalpanā karṇe}
  \item \emph{Say it:} \emph{kalpanā karṇe}, once more
  \item \emph{Say it:} say \emph{kalpanā karṇe}
  \item \emph{Recall:} say the Marathi for to turn, then the Marathi for to enjoy, then say \emph{kalpanā karṇe} again
\end{itemize}

\begin{tcolorbox}[breakable,colback=teal!4,colframe=teal!35!black,title={Before you move on}]
What does \mr{कल्पना} \mr{करणे} mean? (\textquotedblleft{}To imagine\textquotedblright{}.) Say it once more.
\end{tcolorbox}
\section[\texorpdfstring{\mr{कळवणे} (kaḷavṇe)}{कळवणे (kaḷavṇe)}]{\mr{कळवणे} (kaḷavṇe) — to inform, to let know}

\label{lesson:MR-C253-kalavne}



\begin{quote}
Before the new one: say the Marathi for to enjoy, then the Marathi for to imagine.
\end{quote}

\subsection*{You'll want to know: \mr{कळवणे}}
\textbf{\mr{कळवणे}} — \emph{kaḷavṇe} — \textquotedblleft{}to inform, to let know\textquotedblright{}.

\begin{grammarlens}[title={What you've built}]
One more word for this chapter.
\end{grammarlens}

\subsection*{Your turn}
\begin{itemize}
\raggedright
  \item \emph{Say it:} \emph{kaḷavṇe}
  \item \emph{Say it:} \emph{kaḷavṇe}, once more
  \item \emph{Say it:} say \emph{kaḷavṇe}
  \item \emph{Recall:} say the Marathi for to enjoy, then the Marathi for to imagine, then say \emph{kaḷavṇe} again
\end{itemize}

\begin{tcolorbox}[breakable,colback=teal!4,colframe=teal!35!black,title={Before you move on}]
What does \mr{कळवणे} mean? (\textquotedblleft{}To inform, to let know\textquotedblright{}.) Say it once more.
\end{tcolorbox}
\section[\texorpdfstring{\mr{तक्रार} \mr{करणे} (takrār karṇe)}{तक्रार करणे (takrār karṇe)}]{\mr{तक्रार} \mr{करणे} (takrār karṇe) — to complain}

\label{lesson:MR-C253-takrar-karne}



\begin{quote}
Before the new one: say the Marathi for to imagine, then the Marathi for to inform.
\end{quote}

\subsection*{You'll want to know: \mr{तक्रार} \mr{करणे}}
\textbf{\mr{तक्रार} \mr{करणे}} — \emph{takrār karṇe} — \textquotedblleft{}to complain\textquotedblright{}.

\begin{grammarlens}[title={What you've built}]
One more word for this chapter.
\end{grammarlens}

\subsection*{Your turn}
\begin{itemize}
\raggedright
  \item \emph{Say it:} \emph{takrār karṇe}
  \item \emph{Say it:} \emph{takrār karṇe}, once more
  \item \emph{Say it:} say \emph{takrār karṇe}
  \item \emph{Recall:} say the Marathi for to imagine, then the Marathi for to inform, then say \emph{takrār karṇe} again
\end{itemize}

\begin{tcolorbox}[breakable,colback=teal!4,colframe=teal!35!black,title={Before you move on}]
What does \mr{तक्रार} \mr{करणे} mean? (\textquotedblleft{}To complain\textquotedblright{}.) Say it once more.
\end{tcolorbox}
